% Pursuing Stacks (cotes 134-2 to 134-8): archive cote 134-5, site page 36 (PDF page 37), typescript page 223.
% typescript. Draft compared: ps4.tex:1735-1810.
% AI-transcribed by Claude Opus 5.5 (claude-opus-5-5, adaptive thinking, effort high), Message Batches API, batch ps2 (msgbatch_01NVwr1Yuo8qJf5rykeGBXPK), 2026-11-12; not yet checked by eye.
\scanpage{134-5/36}
\mspage{223}

Assume now A is a local test category, i.e.\inserted{$A^\wedge$} admits a separating
interval
\[ \overline{I} = (I,\hand{$\delta$}_0,\hand{$\delta$}_1) \]
such that I is aspheric over the final object of $A^\wedge$. Then $p_1^*(\underline{I})$ is of
course a separating interval in $(AxB)^\wedge$, which by cor.1 is aspheric over
the final object. Hence:

\uline{Proposition} 2: If A is a local test category, so is AxB for any B in
(Cat). If A is a test category, then so is AxB for any aspheric B .

The second assertion follows, remembering that a test category is
just an \uline{aspheric} local test category. Using cor. 2, we get :

\uline{Corollary}: If A is a strict test category, and B totally aspheric,
then AxB is a strict test category. In particular, if A and B are
strict test categories, so is their product.

We need only remember that strict test categories are just test
categories which are totally aspheric.

As an illustration, we get the fact that the categories of multi-%
complexes of various kinds, which we can \inserted{even} take mixed (semisimplicial
in some variables, cubical in some others, and hemispherical say in
others still) are strict modelizers(as generally granted), which corres
\hand{-}ponds to the fact that a \xout{xx} finite product of \xout{the} standard semisimplicial
cubical and hemispherical\inserted{(strict)} test categories \hand{$\Delta$} , \hand{$\Box$} and \hand{\unsure{$\bigcirc$}} , is again a
strict test category.

I would like now to dwell a little upon the comparison of ``homotopy
models'', \inserted[typed]{using}\xout{\illegible{1}} respectively \inserted[typed]{\xout{\illegible{1}}}\xout{in} two test categories A and B, namely
working in $A^\wedge$ and $B^\wedge$ respectively. More specifically, we have two
descriptions of Hot (short for Hot($\underline{W}$) here ) , namely as \xout{Hot$_A$}
\[ \mathrm{Hot}_A = \underline{W}_A{}^{-1}A^\wedge \quad\text{and}\quad \mathrm{Hot}_B = \underline{W}_B{}^{-1}B^\wedge \quad , \]
and we want to describe conveniently the ta\inserted{u}tological equivalence between
these two categories ( this equivalence being defined up to unique iso-%
morphism). The most \hand{``}tautological\hand{''} way indeed is to use the basic modeli-%
zer (Cat) and it's localization (Hot) as the int ermediary \xout{\illegible{1}} ,
i.e. using the diagram of equivalences of categories
\diagram{
(1)\qquad
\begin{tikzcd}[column sep=small]
\mathrm{Hot}_A \ar[dr, "\hand{$\approx$}", "\overline{\hand{$i$}}_{\hand{$A$}}"'] & & \mathrm{Hot}_B \ar[dl, "\overline{\hand{$i$}}_{\hand{$B$}}", "\hand{$\approx$}"'] \\
& (\mathrm{Hot}) &
\end{tikzcd}
\qquad .
}
Remember we have a handy quasi-inverse \hand{$\overline{j}_B$}

\margin{36}
